Retrieval-Augmented Generation (RAG) has emerged as a dominant paradigm for grounding large language model outputs in factual, domain-specific knowledge. 
Yet practitioners face a vast, poorly understood configuration space spanning chunking strategies, embedding models, retrieval methods, and generation models. 
We present a modular, extensible benchmarking framework designed to attach non-invasively to existing RAG pipelines and automatically evaluate pipeline configurations. 
By covering the full RAG lifecycle and automatically exploring a large configuration space across components such as chunking, embedding, retrieval, and generation, the framework eliminates manual tuning while ensuring complete experiment reproducibility via MLflow.
To maximize evaluation coverage, it combines LLM-based quality metrics (RAGAS) with established numerical metrics from information retrieval and natural language generation. 
%The framework is demonstrated on Wikipedia QA datasets, benchmarking over 100 configuration combinations across locally-served open-weight models. 
%The results reveal non-obvious interactions between chunking granularity, retrieval strategy, and generation quality. 
%The framework supports extensible dataset adapters, per-role endpoint configuration for distributed setups, and integrated experiment tracking via Mlflow.

Multiple pre-configured dataset adapters are supported out of the box, enabling evaluation across diverse domains, with the option to integrate custom domain-specific data.
The framework is demonstrated across multiple QA datasets, benchmarking over 100 configuration combinations on locally-served open-weight models.
Results reveal interactions between chunking granularity, retrieval strategy, and generation quality, and demonstrate that optimal configurations differ across systems underscoring the practical value of automated, non-invasive benchmarking that can be plugged into any existing RAG deployment.